\expandafter\def\csname HistoricalNoteText21\endcsname{Die Elemente III ergeben für die Rektaszension am 4. Febr.\ einen um $13'$ und am 28. Juni einen um $32'$ kleineren Wert, für die Deklination einen um $2'$ größeren.}
\SourceLeaf{228}%
\bigskip\Needspace{10\baselineskip}
\pdfbookmark[1]{Gauß an Bode. 1802 Sept. 10}{bode18020910}
\begin{center}GAUSS AN BODE.\end{center}
\begin{flushright}Braunschweig 1802 Sept.\ 10.\end{flushright}
\begin{center}
\medskip\hrule\medskip
Astronomisches Jahrbuch für 1805. Berlin 1802. Seite 227.\par
\medskip\hrule
\end{center}\medskip

Die Elemente, nach denen die folgende Ephemeride für den 9. Hauptplaneten \textit{Pallas Olbersiana} berechnet wurde, beruhen auf den ersten Meridianbeobachtungen vom Seeberg, den letzten aus Greenwich und den Beobachtungen von Dr.\ \textsc{Olbers} vom 8. Juli. Sie lauten wie folgt.
\begin{center}
{\small\setlength{\tabcolsep}{3pt}
\begin{tabular}{l r}
Epoche 1802, 31. März, Seeberger Meridian $\ .\ .\ .\ .\ .$&$162^0\ 55'\quad6''8$\\
Tägliche tropische Bewegung $\ .\ .\ .\ .\ .\ .\ .\ .\ .$&$769''7263$\\
Logarithmus der großen Halbachse $\ .\ .\ .\ .\ .\ .$&$0.4424992$\\
Sonnenferne $\ .\ .\ .\ .\ .\ .\ .\ .\ .\ .\ .\ .\ .\ .\ .$&$301^0\ 38'\ 42''$\\
Exzentrizität $\ .\ .\ .\ .\ .\ .\ .\ .\ .\ .\ .\ .\ .\ .\ .$&$0.243888$\\
Knoten {\AstroSymbols ☊} $\ .\ .\ .\ .\ .\ .\ .\ .\ .\ .\ .\ .\ .\ .\ .\ .$&$172^0\ 26'\ 31''$\\
Neigung $\ .\ .\ .\ .\ .\ .\ .\ .\ .\ .\ .\ .\ .\ .\ .\ .\ .$&$34\ 36\ 59$
\end{tabular}}
\end{center}
Wenn wir die Lichtstärke des Planeten am 8. Juli dieses Jahres $=1$ setzen, ergeben sich folgende Werte.
\begin{center}
\begin{tabular}{r l r}
am 4. April&$\ .\ .\ .\ .\ .\ .\ .\ .\ .\ .$&$=4.23$\\
4. Mai&$\ .\ .\ .\ .\ .\ .\ .\ .\ .\ .$&$=2.75$\\
26. Juli&$\ .\ .\ .\ .\ .\ .\ .\ .\ .\ .$&$=0.80$\\
24. August&$\ .\ .\ .\ .\ .\ .\ .\ .\ .$&$=0.60$\\
\multicolumn{2}{l}{größte Lichtstärke im nächsten Jahr}&$=0.653$
\end{tabular}
\end{center}

Das zeigt, dass man sich noch nicht ganz darauf verlassen kann, den Planeten 1803 wiederzufinden. Die große Höhe über dem Horizont und der größere Abstand von der Dämmerung müssen den Verlust ausgleichen. Doch ich schmeichle mir mit der Hoffnung, dass man den Planeten in südlichen Klimazonen leichter finden wird, etwa in Palermo, wo die Juninächte um Mitternacht ganz dunkel sind.

Wenn ich künftig auch eine Ephemeride der $\CeresSymbol$ nach verbesserten Elementen berechne, werde ich sie Ew.\ — ebenfalls zusenden\HistoricalNote{21}. [Die Ephemeride reicht vom 4. Febr.\ bis zum 28. Juni, jeweils um Mitternacht im Seeberger Meridian, 1803. Sie gibt Rektaszension und Deklination in Abständen von drei Tagen minutengenau an.]
\begin{center}
\begin{tabular}{l}
{\AstroSymbols ☍} 30. Juni. Nachmittag.\\
nach den Elementen III 30. Juni. Vormittag, bürgerliche Zeit.
\end{tabular}
\end{center}
